\SourceLeaf{364}
\begin{center}\bfseries [MUTATIONEN DES RAUMES.]\end{center}
\begin{center}[I.]\end{center}
\begin{center}\bfseries [TRANSFORMATIONEN DES RAUMES.]\end{center}
\begin{center}[1.]\end{center}

Man setze
\[
\begin{aligned}
\BareRoot{\frac{1-p-q-r}{2}}&=a\\
\BareRoot{\frac{1-p+q+r}{2}}&=b\\
\BareRoot{\frac{1+p-q+r}{2}}&=c\\
\BareRoot{\frac{1+p+q-r}{2}}&=d,
\end{aligned}
\]
so ist
\[
\begin{array}{ccc}
p&ad+bc&ac-bd\\
bc-ad&q&ab+cd\\
-ac-bd&cd-ab&r
\end{array}
\]
das allgemeine Transformationsschema der Raumcoordinaten, oder, bedeuten $a,b,c,d$ was sie wollen[, so sind die 9 Grössen des Schemas] proportional folgenden:
\[
\begin{array}{ccc}
\tfrac12(cc+dd-aa-bb)&ad+bc&ac-bd\\
bc-ad&\tfrac12(bb+dd-aa-cc)&ab+cd\\
-ac-bd&cd-ab&\tfrac12(bb+cc-aa-dd).
\end{array}
\]
